\documentclass[11pt]{article}

\usepackage[T1]{fontenc}
\usepackage[utf8]{inputenc}
\usepackage{lmodern}
\usepackage{amsmath,amssymb,amsthm}
\usepackage{booktabs}
\usepackage{geometry}
\usepackage{microtype}
\usepackage[hidelinks]{hyperref}

\geometry{margin=1in}

\newtheorem{theorem}{Theorem}
\newtheorem{proposition}{Proposition}
\newtheorem{corollary}{Corollary}
\newtheorem{remark}{Remark}

\newcommand{\sinc}{\operatorname{sinc}}

\title{\textbf{A Universal Scaling Law for Ellipse Perimeters}}
\author{Wayne Baker\\
\small Independent Researcher}
\date{February 2026; revised July 14, 2026}

\begin{document}

\maketitle

\begin{abstract}
Let
\[
\mathbf q(\theta)=(a\cos\theta,b\sin\theta),
\qquad a\ge b>0,
\]
and let $P_N$ be the perimeter of the inscribed polygon obtained from $N$
equally spaced parameter values.  We prove the exact factorization
\[
P_N
=
\sinc\!\left(\frac{\pi}{N}\right)M_N,
\qquad
\sinc x=\frac{\sin x}{x},
\]
where $M_N$ is the composite midpoint rule applied to the ellipse speed.
Because that speed is periodic and analytic in a strip of half-width
\[
\rho_*=\operatorname{artanh}(b/a),
\]
the midpoint error is exponentially small for every fixed nondegenerate
ellipse.  Consequently,
\[
P_N
=
P\,\sinc\!\left(\frac{\pi}{N}\right)
+
O(e^{-\rho N})
\]
for every $0<\rho<\rho_*$, and hence
\[
P-P_N
=
P\left(
\frac{\pi^2}{6N^2}
-
\frac{\pi^4}{120N^4}
+
\frac{\pi^6}{5040N^6}
-\cdots
\right)
+
O(e^{-\rho N}).
\]
Thus the algebraic relative-error coefficients are universal: they do not
depend on the ellipse axes.  The tripling correction
\[
\widehat P_N
=
P_{3N}+\frac{P_{3N}-P_N}{8}
\]
satisfies
\[
P-\widehat P_N
=
\frac{P\pi^4}{1080N^4}
+
O(N^{-6})
+
O(e^{-\rho N}),
\]
which gives fourth-order convergence after a single fixed-scale
cancellation.  The eccentricity dependence enters through the exponentially
small quadrature remainder; as $b/a\to0$, the analytic strip narrows like
$b/a$, explaining why the asymptotic regime requires $N$ of order $a/b$ up
to a tolerance-dependent logarithmic factor.
\end{abstract}

\section{Introduction}

Polygonal approximation of an ellipse contains two distinct numerical
effects.  First, replacing a smooth arc by its chord introduces a local
chord--arc factor.  Second, sampling the varying speed of the
parametrization introduces a periodic quadrature error.  Separating these
effects yields a particularly simple scaling law.

The central result of this paper is ellipse-specific and exact: the
equal-parameter polygonal perimeter factors into a universal sinc term and a
composite midpoint approximation to the exact perimeter.  Standard
exponential convergence of periodic trapezoidal and midpoint rules then
shows that, for every fixed $b>0$, all algebraic powers of $N^{-1}$ come from
the sinc factor alone.

This statement is narrower and stronger than a blanket claim about arbitrary
chord constructions on arbitrary smooth curves.  It applies to the standard
equal-parameter polygonal approximation of a nondegenerate ellipse.

\section{Exact chord factorization}

Let
\[
h=\frac{2\pi}{N},
\qquad
\theta_j=jh,
\qquad
t_j=\left(j+\frac12\right)h,
\qquad
j=0,\ldots,N-1.
\]
Define the ellipse speed
\[
v(t)
=
\|\mathbf q'(t)\|
=
\sqrt{a^2\sin^2t+b^2\cos^2t}.
\]
The exact perimeter is
\[
P=\int_0^{2\pi}v(t)\,dt
=4a\,E\!\left(1-\frac{b^2}{a^2}\right),
\]
where $E$ is the complete elliptic integral of the second kind in the
parameter convention.

\begin{theorem}[Exact polygon--midpoint identity]\label{thm:factor}
For every integer $N\ge3$,
\[
P_N
=
2\sin\!\left(\frac{\pi}{N}\right)
\sum_{j=0}^{N-1}v(t_j).
\]
Equivalently, if
\[
M_N
=
\frac{2\pi}{N}
\sum_{j=0}^{N-1}v(t_j)
\]
is the composite midpoint rule for $P$, then
\[
P_N
=
\sinc\!\left(\frac{\pi}{N}\right)M_N.
\]
\end{theorem}

\begin{proof}
For one chord, the angle-difference identities give
\[
a\cos(\theta_j+h)-a\cos\theta_j
=
-2a\sin t_j\sin\frac h2
\]
and
\[
b\sin(\theta_j+h)-b\sin\theta_j
=
2b\cos t_j\sin\frac h2.
\]
Therefore its length is
\[
\|\mathbf q(\theta_j+h)-\mathbf q(\theta_j)\|
=
2\sin\frac h2\,
\sqrt{a^2\sin^2t_j+b^2\cos^2t_j}
=
2\sin\frac h2\,v(t_j).
\]
Summing over $j$ proves the first identity.  Since $h=2\pi/N$, the second
follows by multiplying and dividing by $h$.
\end{proof}

\section{Analytic midpoint remainder}

The factorization in Theorem~\ref{thm:factor} reduces the remaining problem
to periodic quadrature.

\begin{proposition}[Analytic strip]\label{prop:strip}
Set $r=b/a$.  For $0<r\le1$, the function $v$ has a $2\pi$-periodic analytic
continuation to every closed strip
\[
|\operatorname{Im}z|\le\rho<\rho_*,
\qquad
\rho_*=\operatorname{artanh}(r).
\]
\end{proposition}

\begin{proof}
The possible branch points occur at zeros of
\[
a^2\sin^2z+b^2\cos^2z
=
b^2+(a^2-b^2)\sin^2z.
\]
For $a>b$, the nearest zeros satisfy
\[
\sin z=\pm i\frac{b}{\sqrt{a^2-b^2}},
\]
whose distance from the real axis is
\[
\operatorname{arsinh}\!\left(
\frac{b}{\sqrt{a^2-b^2}}
\right)
=
\operatorname{artanh}(b/a).
\]
Hence a consistent analytic square-root branch exists in every smaller
strip.  When $a=b$, the speed is constant and the assertion is immediate.
\end{proof}

\begin{proposition}[Periodic midpoint error]\label{prop:midpoint}
For each fixed ellipse with $a\ge b>0$ and each
$0<\rho<\rho_*$,
\[
M_N-P=O(e^{-\rho N}).
\]
\end{proposition}

\begin{proof}
The composite midpoint rule on a full period is a shifted periodic
trapezoidal rule.  For a periodic function analytic and bounded in a strip,
its Fourier coefficients decay exponentially, and aliasing of those
coefficients gives an exponentially small quadrature error.  Applying this
standard result to Proposition~\ref{prop:strip} proves the claim; see, for
example, Trefethen and Weideman~\cite{TrefethenWeideman2014}.
\end{proof}

\section{Universal algebraic scaling law}

\begin{theorem}[Ellipse perimeter scaling law]\label{thm:main}
For any fixed ellipse $a\ge b>0$ and any $0<\rho<\rho_*$,
\[
P_N
=
P\,\sinc\!\left(\frac{\pi}{N}\right)
+
O(e^{-\rho N}).
\]
Consequently, for every fixed integer $m\ge1$,
\[
P_N
=
P\sum_{\ell=0}^{m}
\frac{(-1)^\ell\pi^{2\ell}}
{(2\ell+1)!\,N^{2\ell}}
+
O(N^{-2m-2})
+
O(e^{-\rho N}).
\]
In particular,
\[
P-P_N
=
P\left(
\frac{\pi^2}{6N^2}
-
\frac{\pi^4}{120N^4}
+
\frac{\pi^6}{5040N^6}
\right)
+
O(N^{-8})
+
O(e^{-\rho N}).
\]
\end{theorem}

\begin{proof}
Combine Theorem~\ref{thm:factor} with
Proposition~\ref{prop:midpoint}:
\[
P_N
=
\sinc\!\left(\frac{\pi}{N}\right)
\bigl(P+O(e^{-\rho N})\bigr).
\]
The sinc factor is bounded and has the Taylor series
\[
\sinc x
=
\sum_{\ell=0}^{\infty}
\frac{(-1)^\ell x^{2\ell}}{(2\ell+1)!}.
\]
Substitution gives the stated formulas.
\end{proof}

\begin{remark}
The word ``universal'' refers to the algebraic \emph{relative}-error
coefficients.  For the standard equal-parameter ellipse polygon,
\[
\frac{P-P_N}{P}
\sim
\frac{\pi^2}{6N^2},
\]
independently of $a$ and $b$.  The onset of this asymptotic regime is not
uniform as $b/a\to0$.
\end{remark}

\section{Tripling-based cancellation}

Define the one-step corrected perimeter
\[
\widehat P_N
=
P_{3N}
+
\frac{P_{3N}-P_N}{8}
=
\frac{9P_{3N}-P_N}{8}.
\]

\begin{corollary}[Fourth-order corrected law]\label{cor:corrected}
For every fixed nondegenerate ellipse and every $0<\rho<\rho_*$,
\[
P-\widehat P_N
=
\frac{P\pi^4}{1080N^4}
-
\frac{P\pi^6}{40824N^6}
+
O(N^{-8})
+
O(e^{-\rho N}).
\]
Hence
\[
\frac{P-P_{3N}}{P-P_N}\longrightarrow\frac19,
\qquad
\frac{P-\widehat P_{3N}}{P-\widehat P_N}
\longrightarrow\frac1{81}.
\]
\end{corollary}

\begin{proof}
Apply the expansion in Theorem~\ref{thm:main} at $N$ and $3N$ and form
$(9P_{3N}-P_N)/8$.  The $N^{-2}$ term cancels.  The $N^{-4}$ coefficient is
$-P\pi^4/1080$ in $\widehat P_N-P$, and the next coefficient is
$P\pi^6/40824$.  The ratio limits follow from the leading nonzero terms.
\end{proof}

This is a Richardson-type fixed-scale cancellation.  It does not require
knowledge of the perimeter or of the leading coefficient; only the known
scale factor $3^{-2}$ is used.

\section{Eccentricity and the asymptotic onset}

The algebraic coefficients in Theorem~\ref{thm:main} are independent of
eccentricity, but the exponentially small midpoint remainder is not.  Its
natural strip width is
\[
\rho_*=\operatorname{artanh}(b/a).
\]
For a highly eccentric ellipse,
\[
\operatorname{artanh}(b/a)
\sim
\frac ba
\qquad
(b/a\to0).
\]
Thus suppressing the quadrature remainder to a fixed tolerance requires,
at the level of exponential order,
\[
N
\gtrsim
\frac{\log(1/\varepsilon)}{\operatorname{artanh}(b/a)}
\sim
\frac ab\log\frac1\varepsilon.
\]
This explains an empirical inverse-axis-ratio scaling without asserting a
universal finite-$N$ threshold.  The degenerate case $b=0$ is excluded:
there the speed becomes $a|\sin t|$, analyticity is lost, and the remainder
is no longer exponentially small.

Ignoring the exponential transient, the leading relative-error design
estimates are
\[
N_{\mathrm{raw}}
\approx
\frac{\pi}{\sqrt{6\varepsilon}},
\qquad
N_{\mathrm{corr}}
\approx
\left(
\frac{\pi^4}{1080\varepsilon}
\right)^{1/4}.
\]
They should be used only after $N\operatorname{artanh}(b/a)$ is sufficiently
large.

\section{Numerical verification}

The accompanying script evaluates the exact perimeter using the complete
elliptic integral and computes the polygonal sums directly at high precision.
Table~\ref{tab:numerics} displays the normalized quantities predicted by the
theory.

\begin{table}[ht]
\centering
\small
\begin{tabular}{rrrrr}
\toprule
$b/a$ & $N$ & $N\rho_*$ &
$N^2(P-P_N)/P$ &
$N^4(P-\widehat P_N)/P$ \\
\midrule
$0.90$ & $162$   & $238.46$ & $1.6449031365$ & $0.0901927055$ \\
$0.50$ & $486$   & $266.98$ & $1.6449306301$ & $0.0901935031$ \\
$0.10$ & $1458$  & $146.29$ & $1.6449336850$ & $0.0901935917$ \\
$0.01$ & $13122$ & $131.22$ & $1.6449340621$ & $0.0901936027$ \\
\midrule
Limit & -- & -- & $\pi^2/6=1.6449340668$ &
$\pi^4/1080=0.0901936028$ \\
\bottomrule
\end{tabular}
\caption{Scaled relative errors for four axis ratios.  The differing values
of $N$ place each case well inside its analytic asymptotic regime.}
\label{tab:numerics}
\end{table}

The data confirm two separate conclusions: the algebraic limiting constants
are universal, while increasingly eccentric ellipses require larger $N$ to
make the exponentially small quadrature remainder negligible.

\section{Scope and relationship to geometric N-series refinement}

This theorem supplies an exact ellipse-specific foundation for the
even-power assumption used in fixed-scale N-series acceleration.  It is
compatible with the general scaling-cancellation principle, but it does not
assert that every chord-based construction on every smooth curve has the
same coefficients.

The result is also distinct from nonlinear geometric residual iteration.
Here the active variable is the resolution scale $N^{-1}$, and a linear
combination of fixed-resolution approximants raises the algebraic order from
two to four.

\section{Limitations}

The conclusions apply to:

\begin{itemize}
\item nondegenerate ellipses with $a\ge b>0$;
\item the standard polygon obtained from equal increments of the ellipse
parameter;
\item the asymptotic regime $N\to\infty$ for fixed axis ratio.
\end{itemize}

They do not claim:

\begin{itemize}
\item uniform exponential constants as $b/a\to0$;
\item the same sinc factorization for arbitrary parameterizations or meshes;
\item numerical superiority over direct elliptic-integral algorithms;
\item a universal theorem for all smooth closed curves without additional
hypotheses.
\end{itemize}

\section{Conclusion}

The standard equal-parameter polygonal approximation of an ellipse satisfies
an exact factorization:
\[
P_N
=
\sinc\!\left(\frac{\pi}{N}\right)M_N.
\]
For every fixed nondegenerate ellipse, periodic analytic quadrature makes
$M_N-P$ exponentially small.  The full algebraic error hierarchy therefore
comes from the universal sinc factor, giving relative coefficients that are
independent of the ellipse axes.

A single tripling-based cancellation removes the inverse-square term and
produces fourth-order accuracy:
\[
P-\widehat P_N
\sim
\frac{P\pi^4}{1080N^4}.
\]
Eccentricity controls not the algebraic coefficients but the onset of the
asymptotic regime through the strip width
$\operatorname{artanh}(b/a)$.  This separation gives a rigorous and
transparent scaling law for ellipse-perimeter N-series refinement.

\begin{thebibliography}{9}

\bibitem{Richardson1927}
L.~F. Richardson,
``The Deferred Approach to the Limit,''
\emph{Philosophical Transactions of the Royal Society of London, Series A},
226 (1927), 299--361.

\bibitem{TrefethenWeideman2014}
L.~N. Trefethen and J.~A.~C. Weideman,
``The Exponentially Convergent Trapezoidal Rule,''
\emph{SIAM Review} 56 (2014), no.~3, 385--458.

\bibitem{Borwein1987}
J.~M. Borwein and P.~B. Borwein,
\emph{Pi and the AGM: A Study in Analytic Number Theory and Computational
Complexity},
Wiley, 1987.

\bibitem{BakerScaling2026}
W.~Baker,
\emph{A Scaling--Cancellation Principle for Chord-Based Geometric
Approximations},
preprint, 2026.

\end{thebibliography}

\end{document}
